% ECONOMICS_PROBLEM: {"id": "D71-647", "title": "Core nonemptiness for approval-count committee elections: resolved historical question", "jel_code": "D71", "source_id": "OP-0202"}
% !TeX program = xelatex
\documentclass[11pt,a4paper]{article}
\usepackage[margin=23mm]{geometry}
\usepackage{fontspec}
\setmainfont{Latin Modern Roman}
\usepackage{amsmath,amssymb,mathtools}
\usepackage{xcolor,enumitem,booktabs,longtable,array,xurl,hyperref}
\definecolor{muted}{HTML}{53636C}
\hypersetup{colorlinks=true,urlcolor=blue,linkcolor=blue}
\setlength{\parindent}{0pt}
\setlength{\parskip}{5pt}
\newcommand{\R}{\mathbb R}
\newcommand{\E}{\mathbb E}
\newcommand{\N}{\mathbb N}
\newcommand{\ind}{\mathbf 1}
\newcommand{\Var}{\operatorname{Var}}
\newcommand{\Cov}{\operatorname{Cov}}
\newcommand{\supp}{\operatorname{supp}}
\DeclareMathOperator*{\argmin}{arg\,min}
\DeclareMathOperator*{\argmax}{arg\,max}
\newcommand{\entrymeta}[3]{{\sffamily\small\color{muted}\textbf{\texttt{#1}}\quad|\quad #2\par #3\par}}
\newcommand{\Problemref}[1]{\texttt{#1}}
\newcommand{\JELref}[2]{\texttt{#2} (JEL \texttt{#1})}
\begin{document}
\section*{Core nonemptiness for approval-count committee elections: resolved historical question}
\textbf{Problem ID:} \texttt{D71-647}\quad \textbf{JEL:} \texttt{D71}\par
% BEGIN ECONOMICS STATEMENT
\label{problem:OP-0202}
\label{audited:ECON-007}

% GAME_THEORY_PROBLEM: {"local_id": "GT-072", "original_number": 72, "kind": "S", "entry_kind": "statement_only", "published": false, "review_required": true, "category": "game_theory"}
\label{game:GT-072}
{\sffamily\small\color{muted}
\noindent\textbf{ID: GT-072.}\quad\textbf{Category:} Approval-based committee voting.
\quad\textbf{Kind: S.}\quad\textbf{Status:} Resolved historical question; September 2026 preprint.
\par}
\subsection*{Definitions and mathematical statement}
Let \(N=\{1,\ldots,n\}\), \(n\geq 1\), let \(C\) be a finite nonempty
candidate set, and let \(1\leq k\leq |C|\) be an integer.
Each voter \(i\) approves a set \(B_i\subseteq C\), and has utility
\[
 u_i(T)=|B_i\cap T|\qquad(T\subseteq C).
\]
A committee \(W\subseteq C\), \(|W|=k\), is core-stable if
\[
 \forall\,\varnothing\ne S\subseteq N\
 \forall\,T\subseteq C\text{ with }|T|\leq \frac{|S|k}{n}\
 \exists i\in S:\quad |B_i\cap W|\geq |B_i\cap T|.
\]
The proportional budget is rational; since \(|T|\) is integral, its
constraint is equivalent to \(|T|\leq\lfloor |S|k/n\rfloor\).

\paragraph{Historical question, resolved in the cited preprint.}
For every \(n,C,k\) and every approval profile \((B_i)_{i\in N}\),
does a core-stable committee \(W\) exist?

\subsection*{Short English statement}
The historical question asks whether every approval election can select a
committee that no proportionally funded coalition can unanimously improve upon.
[S3] reports a positive answer and polynomial-time computation in this domain.

\subsection*{Variants and scope boundaries}
All coalition members must strictly improve to block the committee.
The deviating committee \(T\) need not be disjoint from \(W\), and
need not be unanimously approved. Utilities count approved winners.
They are neither arbitrary cardinal scores nor binary utilities of
an entire committee. The case \(k=0\) is trivial.
A committee of size at most \(k\) can be padded to exactly \(k\): every
voter's utility weakly increases, so padding cannot create a strict block.
Thus ECON-007's at-most-budget formulation is the same existence question.
Allowing weak improvement for everyone and strict improvement for one changes
the blocking relation. [S3] also treats quotas \(n/(k+1)<q\leq n/k\) and a
strict Droop boundary convention; the weak boundary \(q=n/(k+1)\) is not
silently included. Non-unit costs, fractional committees and general cardinal
scores require their own status evidence.

\subsection*{Source relationship and status}
[S1] explicitly posed general core nonemptiness in December 2025.
[S2] subsequently established partial results while retaining the general
question. [S3], posted on 10 September 2026, proves nonemptiness and supplies
polynomial-time computation for the approval-count, unit-cost, Hare-quota
domain defined above. Its source identity and theorem scope were checked for
this update; this is a preprint resolution, not an independent verification of
every proof step. The old open-status claim is withdrawn, and this entry is
retained for historical traceability rather than counted as an unsolved
conjecture. ECON-007 is consolidated here without adding a duplicate question.
Entry~\ref{game:GT-066} retains the counterexample to the different
general-additive strengthening.

\subsection*{References}
\begingroup\footnotesize\raggedright
\noindent[S1] Ratip Emin Berker, Emanuel Tewolde, Vincent Conitzer,
Mingyu Guo, Marijn Heule, and Lirong Xia,
\emph{On the Edge of Core (Non-)Emptiness: An Automated Reasoning
Approach to Approval-Based Multi-Winner Voting},
2025; AAAI 2026; arXiv:2512.16895v1.
\url{https://arxiv.org/html/2512.16895v1}.
Locators: Section 2, Definition 2 (core stability);
the introduction (general nonemptiness question).

\noindent[S2] Patrick Becker, Matthias Greger, and Dominik Peters,
\emph{Core Existence in Approval-Based Committee Elections with up to
Seven Voter Types}, 2026, arXiv:2605.06194v2.
\url{https://arxiv.org/html/2605.06194v2}.
Locators: Section 2.1, Definition 2.1; Sections 1 and 9
(partial results and remaining general question).

\noindent[S3] Patrick Becker, Matthias Greger, and Dominik Peters,
\emph{Existence of the Core in Approval-Based Committee Elections},
10 September 2026, arXiv:2609.11912v1; author preprint.
\url{https://arxiv.org/html/2609.11912v1}.
Locators: Section 3 (blocking definition), Theorem 6.1 and Section 7
(existence and computation). Resolving source for the historical main question;
also supplied as [S2] in ECON-007.
\par\endgroup

\subsection*{Common conventions: Source mathematical conventions}

\textbf{Finite games.} For a finite set $A$, $\Delta(A)$ is its probability
simplex. Players choose independent mixed strategies in Nash equilibrium;
correlation is allowed only when explicitly part of the solution concept.
In a game with payoff functions $u_i$ and strategy profile $x$, an additive
$\varepsilon$ Nash equilibrium satisfies
\[
 u_i(x)\geq u_i(x_i',x_{-i})-\varepsilon
 \quad\text{for every player $i$ and every admissible deviation $x_i'$}.
\]
For numerical approximation, payoffs are normalized as locally specified.
An $\varepsilon$ payoff residual is distinct from distance at most
$\varepsilon$ to the set of exact equilibria. Point convergence, convergence
of empirical play, and convergence in distance to an equilibrium set are
also distinct.

\textbf{Computational representations.} Unless stated otherwise, a complexity
question takes rational data with binary encoding; polynomial time is measured
in total bit length. A fixed-accuracy polynomial algorithm is not necessarily
polynomial in $1/\varepsilon$, and neither is necessarily polynomial in
$\log(1/\varepsilon)$. Succinct games and value, demand, or sampling oracles
must declare their interfaces. Exact real solutions require an explicit
representation; an unspecified list of arbitrary real numbers is not a
finite computational output. A claimed algorithmic barrier is meaningful
only for its specified model and reductions.

\textbf{Dynamic games.} Histories, information, observation, and allowable
randomization are part of the model. A stationary strategy depends only on
the current state; a positional strategy in a deterministic graph game
chooses one action at each controlled vertex. Nash equilibrium at an initial
state and equilibrium after every history are different requirements.
An expectation of a pathwise $\liminf$ payoff, a $\liminf$ of expected averages,
a limiting expected average, and a uniform finite-horizon equilibrium are
different criteria. None is silently substituted for another.

\textbf{Allocation and mechanisms.} A complete allocation assigns every item
exactly once unless otherwise stated. Zero-valued goods, negative values,
capacity constraints, feasibility, lotteries, and copies of goods affect
fairness definitions and are specified locally. Dominant-strategy incentive
compatibility, Bayesian incentive compatibility, universal truthfulness,
and truthfulness in expectation impose different inequalities. Whenever
an objective ratio has zero denominator, its convention is stated or those
instances are excluded.

\textbf{Infinite-dimensional models.} Expectations, integrals, stochastic
equations, and optimization objectives require their local regularity,
integrability, filtrations, and admissible domains. A backward--forward
mean-field system is a boundary-value problem; it is not an initial-value
semiflow without an additional construction. Long-time, large-population,
and vanishing-noise limits require an order of limits and a convergence
topology. If a minimum need not be attained, the objective is its infimum
or supremum and the approximation requirement is stated separately.

\subsection*{Common conventions: Additional-source status and mathematical conventions}

\label{audited:conventions}

Of the 100 supplied ECON entries, 65 distinct problems appear here; 32 close
matches refine earlier entries, and three duplicate questions map to existing
entries. Appendix~\ref{app:audited-crosswalk} lists every destination.
The attachment's P/R/S/U distinctions and source audits are retained. Its
precise mathematical formalizations are editorial unless the individual audit
says otherwise. Candidate subsidy targets ECON-004--006 retain their U flags;
the resolved approval-core question ECON-007 updates OP-0202 above.
The following supplied conventions govern both this part and the attributed
ECON refinements elsewhere in the catalogue, subject to local restrictions.

\textbf{Parameterized questions.} A problem family lists inputs that must be fixed before an instance is evaluated: population, horizon, policies, information, data law, model class and welfare weights. These are required choices, not quantities the citations are assumed to specify. Undefined applications should not be treated as identified estimands.

\textbf{Indices and integrability.} Sets of agents, firms and regions are finite unless a population measure is explicitly used. \(i,j\) index units, \(r\) regions, \(t\) dates and \(h\) horizons. \(\mathbb E\) and \(\Pr\) refer to the declared population and law; \(\mathbf1\{A\}\) is an indicator. Every displayed expectation and discounted sum needs existence in the stated domain. Infinite horizons use a discount in \((0,1)\) and a convergence condition; finite horizons may use discount one.

\textbf{Optimization and existence.} A displayed \(\max\), \(\min\), or \(\arg\max\) is conditional on attainment. For finite feasible menus this is immediate; general spaces need continuity and compactness/coercivity or another existence argument. Without attainment, the target is the supremum/infimum and arbitrarily near-optimal feasible choices. \(\inf\varnothing=+\infty\). Empty feasibility and an unidentified target are different issues.

\textbf{Potential outcomes.} \(Y_i(z)\) is the outcome under a specified intervention, not the observed conditional mean among people choosing \(z\). In binary comparisons \(\Delta Y_i=Y_i(1)-Y_i(0)\). Specify assignment, treatment versions, compliance, information and observation horizon. Interference is allowed under a full policy vector; compressed exposure notation needs a justified mapping. No interference, consistency and overlap are substantive assumptions, not automatic consequences of notation.

\textbf{Populations and observation.} Fix baseline eligibility and population weights. Conditioning on post-policy employment, survival, relocation or program uptake can change the population being compared. Death is not ordinary loss to follow-up; outcomes truncated by death need a composite convention, joint survival target, or explicitly defined survivor stratum. Fertility-changing interventions need a population functional rather than an unexplained fixed descendant cohort. Measurement and missingness models must be supplied.

\textbf{Identification.} A structural model \(m\in\mathcal M\) implies observable law \(P_m\) and target \(\theta(m)\). Identification means
\[
 P_{m_1}=P_{m_2}\ \Longrightarrow\ \theta(m_1)=\theta(m_2)
 \quad\text{for all }m_1,m_2\in\mathcal M .
\]
Without identification, the target is
\[
 \Theta(P;\mathcal M)=\{\theta(m):m\in\mathcal M,\ P_m=P\}.
\]
Writing \(\theta(P)\) before establishing identification can hide incompatible latent models. Confidence sets additionally need a sampling law, confidence level, asymptotic sequence and uniformity class. Point identification, coverage and informative precision are separate properties.

\textbf{Mediation.} Total effects of randomized assignment are distinct from cross-world natural indirect effects. Belief/effort-channel effects require a meaningful mediator intervention and additional measurement, exclusion or mediation assumptions. Randomization of the initial policy alone does not supply those assumptions. Report bounds or interventional alternatives when the stronger target is unsupported.

\textbf{Equilibrium.} A counterfactual \(W(z)\), \(p(z)\), or \(\mu_t^z\) requires individual optimization, feasibility, government/resource budgets, market clearing and state-transition laws. State equilibrium existence and selection, or report ranges over compatible equilibria. Initial-state integration includes subsequent endogenous transitions; current-state integration must use the policy-dependent distribution. Reduced-form invariance after policy is not assumed.

\textbf{Welfare accounts.} Scores, utility, health, dollars and ecological quantities require explicit conversion before addition. Fees, taxes, subsidies, wages and extortion payments are transfers between parties; distributional weights can make them welfare relevant, but their face value is not automatically a resource loss. State ownership, geographic boundaries, foreign-income weights and fiscal finance. Do not add profits to household welfare already including dividends, or count land capitalization and the underlying benefit twice. A benefit-cost expression must distinguish gross benefits from net welfare and subtract every real cost exactly once.

\textbf{Derivatives and risk.} Log derivatives require positive arguments; local responses require admissible perturbations and specified equilibrium branches. Differentiation under expectations needs regularity/domination, and boundaries may allow only one-sided derivatives. Risk uses a specified law; ambiguity uses a declared nonempty law set on a common history space. Adaptive policies must be nonanticipative. A robust criterion is an editorial choice unless attributed explicitly.

\textbf{Scope overlaps.} Allocation, collective choice, contracts, household bargaining and coordinated policy overlap economics with optimization and game theory. They are retained because they appeared in the original catalogue; no claim of a strictly game-theory-free selection is made.
% END ECONOMICS STATEMENT
\end{document}
